\section{Complete implementations and measured complexity}
\label{sec:implementation}

\subsection{Matching strict implementations}

The directory \file{strict/} contains three standalone Python modules.
Each imports only \texttt{json} and \texttt{sys}; none imports a proof
producer, axiom database, normalization oracle, or test utility. The public
input is a JSON object with exactly \texttt{goal} and \texttt{proof} fields.
The goal must be closed and satisfy the entire stated grammar. A successful
run exits with status zero and prints exactly \texttt{ACCEPT}; every failed
check rejects.

The three modules use the same proof-tree format, the same context handling,
the same de Bruijn convention, the same logical rules, and the same public
driver. The scalar control has explicit successor and initialized iteration.
It uses crossed eliminations and witness-directed introduction, so it is an
equivalent matching presentation of the scalar theory, rather than a literal
transcription of the uncrossed \texttt{MPrf} declaration.
The raw-$R$ control replaces successor by its recursor abbreviation and keeps
both computations. The canonical-$R$ module additionally normalizes every
recursor construction, validates canonical inputs, and deletes the explicit
zero-computation rule. This is a comparison of complete matching
implementations, not a comparison of a fragment against a complete module.

\begin{center}
\small
\begin{tabular}{lrrrr}
\toprule
Complete module & Tokens & AST & Bytes & Gzip 9\\
\midrule
Scalar $I$, seven rules & 1433 & 1595 & 4923 & 1224\\
Raw $R$, seven rules & 1378 & 1532 & 4753 & 1209\\
Canonical $R$, six rules & 1338 & 1485 & 4611 & 1205\\
\bottomrule
\end{tabular}
\par\smallskip
Matching strict implementations; normalization and both interfaces are charged.
\end{center}


The canonical program saves 312 normalized bytes against the scalar control,
or $6.34\%$, together with 95 significant tokens and 110 AST nodes. Against
the raw-$R$ control it saves 142 normalized bytes. Thus normalization has
been charged, and the resulting strict implementation is still smaller in
this matched design. These programs are intentionally transparent; their
size does not improve the repository's much more compressed arity codec.

Measurements use CPython 3.12.14. Actual docstrings are removed, the complete
module is normalized with \texttt{ast.unparse}, and one final LF is appended.
Significant tokens exclude comments, indentation, dedentation, and line-end
and end-marker tokens. AST nodes are counted with \texttt{ast.walk}.
The last column is gzip level 9 of the normalized module with timestamp zero;
it is an informational coordinate, not the optimization objective. Raw
source hashes and sizes are recorded separately.

\subsection{Checking syntax, normalization, and binders}

In the strict format, $0$ is JSON integer zero, positive integers are
de Bruijn variable indices, and $[n,a,h]$ is an $R$ term whose body binds
index one. Formula tags are \texttt{=}, \texttt{Q}, and \texttt{D}. In a
$D$ matrix, the universal variable has index one and the existential
variable index two. Booleans, floats, negative indices, unknown tags, and
extra fields are rejected where the specification requires an integer,
formula, or fixed rule record.

The smart constructor is the complete normalization operation at a newly
constructed recursor node:
\[
 \operatorname{rec}(n,a,h)=[n,a,\;0\text{ if }n=0\text{ else }h].
\]
Its children have already been normalized by recursive substitution.
Counter and initializer are traversed at the current binding depth; the body
is traversed one binder deeper. Substituted images are shifted under binders.
The implementation therefore normalizes a zero counter newly exposed by
substitution, and it never mistakes a bound state for a free parameter.

The strict test suite contains 86 named certificates across the three
calculi. All 172 command-line replays of the readable and normalized modules
match their expectations. The canonical calculus has 12 positive and 18
negative certificates. Its positive cases include each inference mechanism,
an actually used induction hypothesis, a paired elimination with a genuine
witness and shifted retained context, and zero-counter normalization exposed
by transport. Negative controls include a future witness, an escaping
eigenparameter, an invalid source, a wrong closure base, a compound transport
template, malformed syntax, and use of the removed zero rule.

There are also 480 independent numerical evaluations comparing substitution
with environment evaluation, plus direct tests of nesting and capture.
These are finite executable controls. The universal soundness and
completeness argument is the correspondence between the six checked rule
cases and the calculus proved equivalent above; the numerical tests alone
do not establish that metatheorem.

\subsection{Separate compact implementation refinement}
\label{sec:compact-refinement}

The archive also contains a small refinement of the \emph{existing} guarded
scalar calculus. It has no role in the canonical-$R$ equivalence theorem.
The source transformations pair the already checked major and target
comparisons, use a shorter integer/list discriminator, and compose the
existing scalar-zero administrative truth convention with the guarded
checker. A conservative variant changes only the boolean spelling of the
expected major and preserves all JSON executions of its parent.

\begin{center}
\small
\begin{tabular}{lrrrr}
\toprule
Complete module & Tokens & AST & Bytes & Gzip 9\\
\midrule
Frozen compact byte control & 265 & 303 & 522 & 329\\
Refined compact byte profile & 265 & 306 & 516 & 332\\
Frozen compact AST control & 262 & 281 & 606 & 348\\
Refined compact AST profile & 263 & 280 & 602 & 348\\
Matched strict guarded control & 541 & 620 & 1381 & 633\\
Refined strict guarded checker & 542 & 624 & 1376 & 633\\
\bottomrule
\end{tabular}
\par\smallskip
Existing scalar calculus only. Byte and AST improvements use different modules.
\end{center}


The compact byte and AST coordinates belong to different programs. Their
public-language qualification remains the inherited guarded Eq/$D$ contract:
the compact programs do not themselves validate the public grammar. The
directory includes a separately counted strict goal and integer/list
certificate front end, frozen controls, converted fixtures, an exact
description of the changed neutral marker, and a reproducible measurement
and test script. The compact gain is an implementation result, not a new
arithmetic basis and not the size of the canonical-$R$ checker.

The eager paired comparison may inspect previously unused instruction
fields. Its completeness statement is therefore made for the specified
well-scoped integer/list certificate language, and no equality of arbitrary
malformed-JSON behavior is claimed for that variant. Full statements and
the execution evidence are in the compact directory. No failed experimental
checker is supplied as an accepted implementation.

